\section{Introduction}
\label{sec:intro}

Most of what a modern state collects, it collects from wages. If the share of income paid as
wages falls and the share paid as returns to capital rises, the question is whether the tax
system reaches the second well enough to replace what it loses on the first. The question is old,
it has been asked of robots and of offshoring before it was asked of artificial intelligence, and
it is usually answered with a rate: set the capital rate high enough and the revenue comes back.

This paper measures both sides of that proposition and finds the usual answer unavailable, for a
reason that is not about rates at all.

\paragraph{The condition, and why it is a boundary rather than a comparison.} Replacing the
revenue lost when a dollar of wages becomes a dollar of capital income requires the effective
marginal rate on that capital income to be at least the effective labor rate times the share of
the wage that does not come back as wages. Written with an explicit coefficient for how much
taxable domestic capital income a displaced wage dollar actually generates, the condition becomes
a surface rather than a pair of point estimates, and the surface is what the comparison turns on.
Section~\ref{sec:fiscal} builds it.

Measured on fourteen biennial vintages of displaced worker data, the share of a displaced wage
bill that returns as wages is \result{RetainedWageShare}. Assembled from sourced components on a
consistently federal basis, and after decomposing the shareholder layer by holder as
Section~\ref{sec:annual} does, the effective marginal wedge on the income that replaces those
wages is \result{CfBaseFed} percent, against the \result{RequiredFedEasy} to \result{RequiredFedHard}
percent replacement requires. A first pass on an undecomposed layer gave \result{CfBaseOldFed} percent,
and that figure is superseded wherever it still appears in the derivation. On that object the condition closes only
where taxable capital income per displaced wage dollar reaches \result{GStarFedEasy} to
\result{GStarFedHard}, and with output preserved that coefficient cannot exceed one.

\paragraph{Which effective rate is used decides the answer.} A marginal wedge is not a rate of
revenue collection. That distinction is standard, and \citet{DevereuxGriffith2003} is its
reference; we claim no part of it, only that the replacement question has been answered without
it. Section~\ref{sec:annual} carries it through the layers. In a stationary state expensing and economic depreciation give the same annual deduction, so
the annual entity base is the whole of net capital income and not rents alone. On that object the
rate is \result{TauAnnFedDec} percent centrally and \result{TauAnnFedLo} to
\result{TauAnnFedHi} percent at every corner of the rate box with pass-through at one, so the condition closes
everywhere, and the rent share that drives the marginal rate's entire spread drops out of the
algebra. The gap between the two objects is \result{AnnualMarginalGapLo} to
\result{AnnualMarginalGapHi} points, more than any disagreement among the components.

We do not treat either rate as the answer, and we do not claim the two bracket a transition
path. Our calculation is stationary, and a stationary comparison cannot establish what is
collected while the capital stock is still building; the direction is suggestive and the
magnitude is not ours to assert. Which rate is the right comparator depends on the horizon over
which displaced wage receipts must be replaced, and on a transition we do not model. Everything
below that was framed in an earlier version of this work as a shortfall is therefore reported
here as conditional on the marginal reading. The rest of the paper is about what is true on both
objects and what is not.

\paragraph{Which lever binds, and the rival explanation.} A gap of that size on the marginal
rate invites the obvious diagnosis, that multinationals book their profits elsewhere. Testing that requires
care about which question is being asked. Moving a parameter across its uncertainty range says
how much the answer would change if we were wrong about it. It does not say how much that
channel costs. A well-measured channel has a narrow range and a small swing, and may still be
the largest thing in the calculation.

So we compare counterfactuals instead, and they are computed on the decomposed layer of
Section~\ref{sec:annual} rather than on the layer an earlier draft used. Eliminating profit
shifting entirely, so that no rent is booked abroad, raises the marginal rate from \result{CfBaseFed}
to \result{CfNoShiftFed} percent federally, a gain of \result{CfGainShiftFed} points. That
\emph{clears} the \result{RequiredFedEasy} percent easier reading while staying below the
\result{RequiredFedHard} percent harder one, and on the all-government basis it
\result{CfShiftClearsAllGov} the easier reading too. An earlier draft said it still fell short, which
was true of the superseded baseline and is not true of this one. Completing the reach of the shareholder tax,
so that every share sits in a taxable account with nothing deferred and nothing stepped up,
raises it to \result{CfOwnReachFed} percent, a gain of \result{CfGainOwnFed} points.

Both figures are smaller than the earlier draft reported, which had a baseline of \result{CfBaseOldFed}
percent and an ownership gain of \result{CfGainOwnOldFed} points. The ownership channel shrinks
because the baseline it is measured against is larger once withholding and the dividend split
are counted: there is less friction left to remove. The ratio of the ownership channel to the
profit-shifting one falls from \result{CfRatioOldFed} to \result{CfRatioFed}, so the ordering survives
while the margin narrows.

Ownership is two mechanisms rather than one, and bundling them would overstate the comparison, so
Section~\ref{sec:levers} separates them by Shapley value over every order of removal. We do not
quote those point attributions here. They were computed on the shareholder layer as first
assembled, and decomposing that layer changes what the mechanisms mean, so the attributions do
not carry over and we have not recomputed them. What survives is the ordering on the marginal
rate, that the taxable shareholder base is the largest of the three and profit shifting the
smallest, and Section~\ref{sec:levers} marks its table accordingly.

Neither counterfactual is a policy. No reform abolishes profit shifting or places all equity in
taxable hands, and the two rows bound what each channel can contribute rather than proposing
anything. The claim is comparative, not dismissive, and it holds on the marginal rate only: on the annual
rate profit shifting moves the answer more than the shareholder base does.

Who owns the claims still governs the shareholder layer, but less of the ownership is beyond
reach than the headline share suggests. About \result{ThetaUntaxable} percent of US corporate
equity sits outside taxable accounts; decomposed by treatment, \result{TradRetireShare} percent is
in traditional retirement arrangements, untaxed on the margin and taxed at ordinary rates on
distribution, \result{ForeignEquityShare} percent is held abroad and bears withholding on dividends,
and only \result{NoOwnerTaxShare} percent bears no owner-level federal tax at all, those
holdings still bearing the entity tax. The five groups sum to the \result{GroupsSum} the source
itself rounds to rather than to one. Of the gains that
accrue in taxable accounts, \result{GainsUntaxedAtDeath} percent are held until death and never
taxed on 2014 data \citep{CBO2014}. The owner's layer is \result{LayerPublished} percent as
first assembled and \result{LayerAnnualDecomp} percent once the treatments are separated. That is a fact about the structure of
ownership. A legislature can change a rate more easily than it can change that, which is a claim
about political feasibility and not about arithmetic: Equation~\ref{eq:condition} is satisfied by
any sufficiently large effective rate, and nothing here shows that no rate could close the gap.

\paragraph{One structure, three places it shows up.} The same ownership concentration appears
in three questions usually kept apart. We claim it appears in all three, not that it explains
them to a common quantitative standard: the revenue side is assembled from sourced components,
the household side is incidence arithmetic, and the transmission side rests on two historical
episodes that contain much besides an equity fall. The third is illustrative and is labelled so
wherever it appears.

\begin{enumerate}[leftmargin=1.6em]
  \item \textbf{The revenue side.} The shareholder layer is thin because most corporate claims
  sit where the individual income tax does not reach currently, though only
  \result{NeverTaxedShare} percent of them escape federal tax altogether.
  Section~\ref{sec:fiscal}.
  \item \textbf{The household side.} Debt-service capacity cannot be restored by broadening
  ownership, because \result{HHEquityOutsideHouseholds} percent of corporate equity sits outside
  the household sector entirely and another \result{HHEquityInRetirement} percent sits in
  vehicles that cannot be spent. Broadened ownership without liquidity moves the boundary by
  exactly zero, which follows from how the exercise defines spendable resources rather than from
  the data, and is reported as an accounting consequence rather than as a finding.
  Section~\ref{sec:household}.
  \item \textbf{The transmission side.} An equity bust transmits weakly to households through the
  wealth channel we measure, because the assets that fall are held by households least likely to
  change their
  spending. Section~\ref{sec:bust}.
\end{enumerate}

\paragraph{Research questions and objectives.}

\begin{description}[leftmargin=2.2em,style=nextline]
  \item[RQ1] What effective rate does the existing tax system reach on income that shifts from
  wages to capital, and what would replacement require? \\
  \textbf{RO1.} Assemble both sides of the condition from sourced components, report it on each
  jurisdiction consistently and on both the marginal and the annual rate, and state it as a
  boundary over the pass-through coefficient. Section~\ref{sec:fiscal}.
  \item[RQ2] Which component of that rate binds, is it profit shifting, and does the answer
  depend on which rate is used? \\
  \textbf{RO2.} Move every component across its whole sourced range with the others at central
  values, rank the levers by measured influence, and report the ranking on both rates.
  Section~\ref{sec:fiscal}.
  \item[RQ3] Does the same constraint appear on the household and transmission sides? \\
  \textbf{RO3.} Measure the incidence of a labor-to-capital shift across the wealth
  distribution, and price the opposite direction through the same accounts.
  Sections~\ref{sec:household} and \ref{sec:bust}.
\end{description}

\paragraph{What this paper does not claim.} We do not forecast automation; the size of any shift
is an input here, never an output. We do not claim the fiscal mechanism, which is established in
the work reviewed in Section~\ref{sec:related}. We do not claim the result settles which tax base
is the right one: we argue for a marginal, flow-based rate and state the case against our own
choice in Section~\ref{sec:fiscal}. And two components of the assembled rate are swept rather
than sourced, so the central figure is labeled scenario wherever it appears and reported with
the sweep around it.
